\subsection{与本项目相关的研究工作积累}
申请人长期从事行为与实验经济学研究，在风险偏好理论、博弈论实验、同行评估行为、田野实验和计量方法等方面已形成系统积累。以下逐项列出申请人的代表性成果，并说明各项成果与本项目的对应关系。

\subsubsection{已发表论文}

[1] Li L, Maniadis Z, Sedikides C. Anchoring in Economics: A Meta-Analysis of Studies on Willingness-To-Pay and Willingness-To-Accept. Journal of Behavioral and Experimental Economics, 2021.

该研究系统检索并编码了经济学领域锚定效应的实证文献，完成了效应量汇总、异质性分析和发表偏倚检验。与本项目的关系： 元分析方法训练直接服务于本项目对国际部分随机资助机制实证证据的系统评估（辅助方法一：文献计量元分析），其中效应量编码与元回归分析框架可迁移至对新西兰HRC、瑞士SNSF等机构实施数据的定量综合。

[2] Boukouras A, Jennings W, Li L, Maniadis Z. Can Biased Polls Distort Electoral Results? Evidence from the Lab. European Journal of Political Economy, 2023.

该研究采用实验室实验方法，检验了有偏信息对选民行为的因果效应，积累了受控实验中信息操纵、被试激励和行为测量的完整设计经验。与本项目的关系： 实验中"有偏信号→行为扭曲"的因果识别范式与本项目"排序压力→评审保守"的核心逻辑一致；其信息操纵设计为本项目实验中评审噪声水平的处理组设计提供了方法参考。

[3] Pantavou K, Fillon A, Li L, Maniadis Z, Nikolopoulos G. Thermal Indices for Assessing the Impact of Outdoor Thermal Environments on Human Health: A Systematic Review of Observational Studies. International Journal of Biometeorology, 2025.

该研究采用PRISMA框架完成了大规模系统综述。与本项目的关系： 进一步巩固了申请人在文献系统检索和规范化综述方面的方法能力，支撑本项目对部分随机资助机制国际实践的证据地图构建。

工作论文（投稿/返修中）

[4] Li L, Louis P, Maniadis Z, Xefteris D. Reciprocity in Peer Assessments. Economic Journal, reject \& resubmit.

该研究构建了基于Levine（1998）互惠偏好理论的正式博弈模型，理论推导了在同时评估与序贯评估两种结构下互惠动机如何导致评估者系统性高估的均衡预测，并通过预注册在线实验（Prolific平台，oTree编程，AEA RCT Registry: AEARCTR-0011281）验证了模型推论。与本项目的关系： （a）该研究直接涉及同行评估中的策略性行为——科研资助评审本质上是同行评估的特殊形式，研究中"评审结构如何塑造评审偏差"的理论框架和实验范式可直接迁移至本项目"排序竞争制vs.合格筛选制"的机制比较（Q1）；（b）博弈论建模能力支撑本项目理论建模方法中评审者与申请者的策略互动分析；（c）在线实验的平台运营、大样本管理和预注册规范流程的实操经验支撑本项目实验室实验的实施。该论文在Economic Journal获得reject \& resubmit，表明研究的学术价值和方法水准已获顶级期刊的基本认可。

[5] Li L, Maniadis Z. A Replication and Robustness Check of the Centipede Game. European Economic Review, revise \& resubmit.

该研究在UCLA CASSEL实验室实施，通过操纵信息反馈结构和收益参数，系统考察了聚合信息对博弈均衡收敛的影响。核心发现：聚合信息反馈能促进纳什均衡收敛，但微小的收益结构变动即可逆转处理效应方向。与本项目的关系： （a）该发现对本项目具有重要方法论启示——在设计资助机制实验时，需特别关注参数敏感性和稳健性检验；（b）实验中关于信息反馈如何影响策略收敛的机制，与本项目中申请者如何响应竞争信号（Q2）密切相关。该论文在European Economic Review获得revise \& resubmit，表明其实验设计和方法论贡献获得一流综合经济学期刊认可。

[6] Li L. Accounting for the Instability of Risk Preferences: Salience Theory versus Cumulative Prospect Theory.

该研究独立设计了包含40道风险选择题的实验室实验（z-Tree编程），在结构估计框架下对显著性理论（Bordalo et al., 2012）与累积前景理论（CPT）进行了系统比较，实现了对概率权重函数、损失厌恶参数和局部思维参数的个体层面估计。与本项目的关系： 前景理论框架下的形式化建模与结构估计能力直接服务于本项目Q1和Q2中评审者损失厌恶参数与申请者风险权重的理论刻画和实验测量；该研究是申请人将行为决策理论应用于具体实验情境的代表性工作。

[7] Li K, Li L, Si W, Xu Z. The Shadow Mommy Effect in Hiring: Evidence from Field and Vignette Experiments.（共同第一作者）

该研究向北京、上海、广州、深圳四个一线城市的真实招聘广告投递了35,713份虚拟简历，结合1,800人的情境实验（vignette experiment），通过"田野实验揭示行为模式+情境实验分离因果机制"的多方法设计，识别了雇主基于生育风险预期的统计性歧视。与本项目的关系： （a）证明申请人具备设计和实施大规模田野实验的全流程能力；（b）"田野揭示行为+实验分离机制"的多方法设计思路与本项目"实验室因果识别+田野DID评估"的方法架构高度一致（Q3）。

[8] Wang H, Li L. Leading by Identification: How Board Chairs' Senior Roles in Industry Associations Shape Environmental Management Practices in Chinese Firms.

该研究使用中国A股上市公司2012—2023年面板数据，运用双向固定效应模型、工具变量法和中介效应分析考察了制度身份对企业环境行为的影响。与本项目的关系： 双向固定效应和工具变量策略为本项目田野/准自然实验部分的双重差分（DID）设计提供了直接的方法经验。

进行中研究

[9] Gall T, Li L, Maniadis Z. Strategic Misreporting in a Committee.（实验设计推进中）

该研究以类似NSFC评审委员会的评分汇总机制为对象，构建理论模型刻画委员会成员在综合评分制度下的策略性误报行为。与本项目的关系： 理论模型直接刻画了评审委员在排序压力下的策略性评分动机，与本项目Q1"排序压力如何塑造评审偏差"高度吻合；实验设计经验可直接服务于本项目实验室实验的参数设定。

[10] Boukouras A, Li L, Maniadis Z, Xefteris D. How Do Biases in Opinion Polls Affect Elections in Costly Voting Environments?（实验设计推进中）

[11] Hanaki N, Li L, Lykopoulos E, Sonsino D. Preference-Identification in Retail Investment Decisions.（实验设计已完成）

[12] Baddeley M, Li L, Zhang J. Naive Herding and House Prices.（实验设计已完成）

[13] Liu M, Li L, Liu P, Wang J. Reward vs Punishment in Collective Action: A Lab-in-the-Field Experiment in China's Pastoral Area.（论文撰写中）

上述进行中研究[10]-[13]虽与本项目主题不直接相关，但体现了申请人在博弈论实验（[10]）、偏好识别实验（[11]）、行为经济学田野实验（[12]-[13]）等方面的广泛方法积累和国际合作网络。

编程与开源能力

申请人熟练掌握Python、R、Stata、JavaScript和Matlab，具备oTree和z-Tree双平台实验编程能力。已在GitHub（github.com/lunzheng-li）公开多个实验程序代码仓库，包括oTree实验程序、数据分析复制包等，具备使用Python/Mesa平台构建智能体仿真模型的技术基础。

对本项目三个科学问题与混合方法设计的整体支撑

综合以上成果，申请人的研究积累对本项目形成如下对应支撑：Q1（评审者行为）——[4]和[9]直接研究同行评估与委员会评分中的策略性偏差，提供理论建模范式和实验测量方案；Q2（申请者行为）——[6]的前景理论结构估计为刻画参照点转移与风险权重变化提供分析工具，[5]的信息反馈实验为理解竞争信号响应提供行为证据；Q3（资助组合结构）——[7]的多方法设计和[8]的面板数据计量提供从微观行为到宏观效应的方法模板。在混合方法层面：理论建模由[4][6][9]支撑，实验室实验由[2][4][5][6]的全链条经验支撑，田野/准自然实验由[7][8]支撑，文献计量由[1][3]支撑，智能体仿真由Python/Mesa编程能力支撑。

\section{研究风险识别与应对措施}

风险一：田野实验合作推进困难。 在国内省/市级资助机构的真实评审流程中试点"合格筛选+部分随机"机制，涉及制度敏感性和组织协调成本。国家机关的评审制度具有较强的路径依赖性，短期内改变现行流程的可能性有限，可能面临合作机构难以落实的风险。

应对措施——三条平行替代路径：

（a）私人资助机构合作。 相比国家资助机构，企业基金会和私人科研资助机构（如腾讯科学探索奖、阿里巴巴达摩院科研基金等）在评审流程设计上具有更大的灵活性和创新意愿。申请人将主动联系此类机构，探索在其部分资助批次中嵌入"合格筛选+随机分配"环节，形成真实资助情境的田野实验数据。

（b）塞浦路斯大学校级科研经费分配试点。 申请人与塞浦路斯大学保持持续合作关系。该校的校内科研经费分配涉及院系层面的同行评审过程，规模适中、制度弹性较大。申请人计划与合作方协商，在某些院系的校内科研基金评审中引入部分随机机制作为试点，同期其他院系沿用传统排序制作为对照组，从而构建校级层面的DID自然实验。该方案的优势在于制度成本低、组织协调可控，且具有跨文化比较价值。

（c）学术会议同行评审数据。 学术会议的论文评审本质上也是一种同行评审竞争性选拔过程，与科研资助评审具有高度相似的结构特征。近期Goldberg, Fanti和Shah（2025）已成功利用NeurIPS 2024和ICLR 2025等顶级会议的真实评审数据，验证了部分随机选拔算法（MERIT）的性能。申请人将参照该思路，联系相关学术会议组织者，探索在部分会议的论文接受环节试行"合格筛选+随机抽签"机制，并收集评审数据用于机制效应分析。学术会议的评审周期短、组织结构相对灵活，是介于实验室实验与大规模田野实验之间的理想过渡情境。

（d）最低产出保障。 即使上述三条田野路径均未能在项目周期内落实，实验室实验仍能独立提供对核心命题（P1—P3）的因果识别证据，辅以国际已有政策变动数据（新西兰HRC、瑞士SNSF）的准自然实验分析和智能体仿真的长期动态推演，确保项目的基本学术产出。

风险二：国际数据获取受限。 新西兰HRC、瑞士SNSF等机构的详细资助数据可能受到数据保护政策限制，影响文献计量和准自然实验分析。

应对措施：（a）优先使用Web of Science、Scopus等公开论文数据库中的项目资助信息字段（Funding Information），结合项目编号和年份信息构建准自然实验数据集。（b）通过现有国际合作网络（申请人博士后合作导师Zacharias Maniadis与欧洲多国科研政策研究者有合作关系）申请数据访问权限。（c）若上述渠道均不可行，将使用NSFC公开数据库中的资助数据进行国内科研资助效率分析，替代国际比较分析。

风险三：多方法整合的技术挑战。 理论模型、实验、仿真和计量分析的参数传递和结论对齐可能面临技术复杂性。

应对措施：（a）采用"理论先导→实验校准→仿真推演"的序贯流程，每一环节的输出作为下一环节的输入参数，形成清晰的信息流。（b）在理论建模阶段即明确可检验预测的操作化定义，确保实验因变量与理论变量的一一对应。（c）在实验设计阶段预留足够的处理组和样本量，使用功效分析（power analysis）确定最小可检测效应量，降低实验结果与理论预测不一致时的解释困难。

风险四：实验结果与理论预测不符。 行为实验的结果可能偏离理论模型的均衡预测。

应对措施：理论建模阶段将构建多个模型变体（含不同行为假设），生成差异化预测，使实验结果即使不支持某一模型也能为模型选择和理论修正提供证据。此外，实验设计中将同时采集多个中介变量和调节变量，确保即使核心效应量低于预期，仍可从机制层面提供有价值的学术贡献。
